\section{Introduction}
\label{sec:intro}

% Problem statement
Autonomous marine vessels have the potential to revolutionize environmental
monitoring, oceanographic data collection, and maritime logistics by
eliminating the need for crewed expeditions. However, existing solutions are
often prohibitively expensive, reliant on motor-driven propulsion with limited
range, or lack the robustness required for extended blue-water operation. A
low-cost, sail-powered autonomous vessel would enable persistent ocean
presence for data collection without the fuel constraints and maintenance
overhead of motorized platforms.

% Solution space
Several approaches to autonomous sailing have been explored in both academic
and commercial domains. The SailBot competition community has produced
numerous small-scale autonomous sailboats, typically using cloth sails with
servo-actuated sheets. Commercial efforts such as Saildrone~\cite{saildrone}
employ rigid wingsails on large platforms for ocean
research~\cite{meinig2019}. University projects have explored various hull
configurations, sail types, and control algorithms ranging from simple
reactive controllers to machine learning approaches.

% Related work summary
Prior autonomous sailboat projects have demonstrated the viability of viability of Global Positioning System (GPS)
waypoint navigation combined with wind-aware path planning. The
self-trimming wingsail concept, where a tail vane provides passive torque to
rotate the main sail element, has been shown to minimize power consumption
compared to motor-actuated sail systems.

% Solution and contribution
ALAN is a semi-autonomous sailboat designed to navigate between
user-defined waypoints without human intervention. The vessel employs a
self-trimming rigid wingsail for propulsion, minimizing power draw while
maximizing aerodynamic efficiency. A dual-core STM32H755 microcontroller
running FreeRTOS manages real-time sensor fusion of the GPS receiver,
inertial measurement unit (IMU), magnetometer, and wind vane, and runs the
proportional-integral-derivative (PID) control loops for rudder and sail
actuation. Long-range communication is achieved via LoRa radio at \SI{915}{\mega\hertz}, providing
telemetry uplink and command and route downlink from a ground station.

% Scope
The scope of this project included:
\begin{itemize}
\item hull design and fabrication;
\item keel and rudder construction;
\item a rigid wingsail with a self-trimming mechanism;
\item sensor integration for navigation and orientation;
\item long-range wireless communication;
\item embedded firmware for autonomous waypoint navigation.
\end{itemize}
Solar charging, full blue-water endurance testing, additional environmental
sensors, and machine-learning-based control were treated as stretch goals
and excluded from scope. The demonstration plan was autonomous navigation
between a series of waypoints on open water, recorded beforehand and shown by
video at Demo Day.

This paper describes the vessel as built. Section~\ref{sec:background}
summarizes the sailing principles that drove the design, and
Section~\ref{sec:related} reviews related work. Section~\ref{sec:overview}
gives the system requirements and architecture. Sections~\ref{sec:mech},
\ref{sec:elec}, and~\ref{sec:sw} describe the mechanical, electrical, and
software subsystems, including the changes made since the proposal.
Section~\ref{sec:test} reports testing, simulation, and the final
demonstration. Section~\ref{sec:impact} discusses safety, societal, and
ethical considerations, Section~\ref{sec:risks} reviews the project's risks
and how they were handled, and Section~\ref{sec:mgmt} covers the team,
schedule, and cost. Section~\ref{sec:conclusion} concludes, and the appendix
lists the bill of materials.